% =====================================================================
%  MarkItDocs — Starter "Plan de trabajo" (licencia MIT)
%  Plan de proyecto: objetivos, fases, hitos, riesgos y equipo.
%  EDITA el contenido de ejemplo.
% =====================================================================
\documentclass[11pt,a4paper]{article}

\usepackage{iftex}
\ifPDFTeX
  \usepackage[utf8]{inputenc}
  \usepackage[T1]{fontenc}
  \usepackage{lmodern}
\else
  \usepackage{fontspec}
\fi
\usepackage[spanish,es-noshorthands]{babel}
\usepackage[top=2.3cm,bottom=2.5cm,left=2.3cm,right=2.3cm]{geometry}
\usepackage{xcolor}
\usepackage{booktabs}
\usepackage{longtable}
\usepackage{array}
\usepackage{titlesec}
\usepackage[hidelinks]{hyperref}

\definecolor{marca}{HTML}{B08D57}
\setlength{\parindent}{0pt}
\setlength{\parskip}{0.5em}
\renewcommand{\arraystretch}{1.3}
\titleformat{\section}{\large\bfseries\sffamily\color{marca!60!black}}{\thesection.}{0.5em}{}

\begin{document}

\begin{center}
  {\LARGE\bfseries Plan de trabajo — [Nombre del proyecto]}\\[4pt]
  {\color{marca}\rule{0.4\linewidth}{2pt}}\\[6pt]
  {\small Responsable: [nombre] · Fecha: \today · Versión: 1.0}
\end{center}

\section{Objetivos}
\begin{itemize}
  \item Objetivo general: qué se quiere lograr al final del proyecto.
  \item Objetivo específico 1: medible y con fecha.
  \item Objetivo específico 2: medible y con fecha.
\end{itemize}

\section{Fases y actividades}
\begin{longtable}{p{2.6cm} p{6.2cm} p{3.4cm} p{2.2cm}}
  \toprule
  \textbf{Fase} & \textbf{Actividades} & \textbf{Entregable} & \textbf{Duración}\\
  \midrule
  1. Análisis    & Levantamiento de requisitos, referencias y contenidos & Documento de alcance & 1 semana\\
  2. Diseño      & Wireframes, diseño visual, revisión con cliente       & Maquetas aprobadas   & 2 semanas\\
  3. Desarrollo  & Implementación frontend/backend, integraciones        & Sitio en preproducción & 3 semanas\\
  4. Lanzamiento & Pruebas, despliegue, dominio y SSL                    & Sitio en producción  & 1 semana\\
  5. Seguimiento & Soporte post-lanzamiento y ajustes                    & Informe de cierre    & 4 semanas\\
  \bottomrule
\end{longtable}

\section{Hitos}
\begin{longtable}{p{5.2cm} p{3.2cm} p{6cm}}
  \toprule
  \textbf{Hito} & \textbf{Fecha objetivo} & \textbf{Criterio de aceptación}\\
  \midrule
  Alcance aprobado      & [fecha] & Documento firmado por el cliente\\
  Diseño aprobado       & [fecha] & Maquetas validadas (máx. 2 rondas)\\
  Entrega en producción & [fecha] & Sitio publicado y verificado\\
  Cierre del proyecto   & [fecha] & Informe final entregado\\
  \bottomrule
\end{longtable}

\section{Riesgos y mitigación}
\begin{longtable}{p{5.6cm} p{2.2cm} p{6.6cm}}
  \toprule
  \textbf{Riesgo} & \textbf{Impacto} & \textbf{Mitigación}\\
  \midrule
  Retraso en contenidos del cliente & Alto  & Fechas límite acordadas y recordatorios\\
  Cambios de alcance                & Medio & Control de cambios presupuestado aparte\\
  Indisponibilidad del hosting      & Bajo  & Proveedor con SLA y copias de seguridad\\
  \bottomrule
\end{longtable}

\section{Equipo y responsabilidades}
\begin{itemize}
  \item \textbf{[Nombre]} — dirección del proyecto y desarrollo.
  \item \textbf{[Nombre]} — diseño UI/UX.
  \item \textbf{[Cliente]} — contenidos y aprobaciones.
\end{itemize}

\end{document}
